\scanpage{39}
tous les pts de $X^{(*)}$.

\deleted{Pour \unsure{les}}
\deleted{Tout \unsure{On veut donc} que les $X^{(a)}$ se recollent \unsure{$(X,\eta)$ correspondants}: $\forall$ couple $(a,b)$, soit $X^{(a,b)}$ l'ens des $u\in\unsure{X^{a,a}|X^{(a)}}$ qui sont \unsure{apparentés} à un pt de $X^b$. D'après ce qui précède \inserted{et $\underset{\leftarrow}{L}_2$ \unsure{il résulte} que c'est} \unsure{on a} un ens \uline{ouvert}. De plus, on a un}

\uline{Lemme} Soit $(Z,\eta)$ ($Z$ loc. de t. fini sur $S$, $\eta\in F(\unsure{S|Z})$) tel que pour $\forall z\in Z$, $\eta_z$ soit modulaire. Considérons $h_Z\xrightarrow{\eta}F$, \deleted{et}
Alors $h_Z\times_F h_Z\subset h_{Z\times Z}$ est représentable par un sous-préschéma de $Z\times Z$, soit $Z'$. Les projections
\diagram{\begin{tikzcd}
Z'\arrow[r,shift left=.6ex,"p_1"]\arrow[r,shift right=.6ex,"p_2"']&Z
\end{tikzcd}}
sont des isomorphismes locaux, [et font de $Z'$ un \deleted{préschéma} graphe d'équivalence dans $Z$ \unsure{$\ldots$}].
\margin{Reprendre les conditions \inserted{nécessaires} sur $F$. Elles sont: \unsure{bien plus faibles}\ldots}

\unsure{AQT}
Or sous les conditions précédentes, prenant $Z=X^{(*)}$, on voit que quand on se restreint au dessus de $S_a$, alors
\diagram{\begin{tikzcd}
Z'_a\arrow[r,shift left=.6ex]\arrow[r,shift right=.6ex]&Z_a
\end{tikzcd}}
est la relation d'équivalence \uline{triviale}. Donc
\mspage{38}
